\chapter{Methodology}
\label{ch:method}

This chapter describes how the corpus was built, how the labels were
derived, how candidate arrangements are generated, and how they are
scored. The models that consume those candidates are the subject of
Chapter~\ref{ch:models}.

\section{The corpus}
\label{sec:corpus}

Three public floor-plan datasets, each in a different native format, are
converted into one representation: a column arrangement and about twenty
geometric measurements per plan. Table~\ref{tab:corpus} gives the
composition.

\begin{table}[H]
\centering
\caption[Corpus composition]{Corpus composition. Median columns and median
regularity are computed after the labelling procedure of
Section~\ref{sec:labels}. ResPlan carries no per-wall thickness, so it
cannot take that procedure and is retained for descriptive purposes only.}
\label{tab:corpus}
\small
\setlength{\tabcolsep}{4pt}
\begin{tabular}{@{}lrllrrr@{}}
\toprule
 & & & & & Median & Median \\
Source & Plans & Native format & Units drawn & Labels & cols. & reg. \\
\midrule
ResPlan & 16,812 & pickled polygons & $\approx$ decimetres & none & 120 & 0.341 \\
Swiss Dwellings & 8,262 & polygons (CSV) & metres & verified & 40 & 0.466 \\
CubiCasa5K & 0 & SVG drawings & drawing pixels & relative & nan & nan \\
\midrule
\textbf{Total} & \textbf{29,279} & & & & & \\
\bottomrule
\end{tabular}
\end{table}

The three sources are drawn in three different unit systems, and no
reliable conversion exists between them: Swiss Dwellings is in metres,
CubiCasa5K in drawing pixels of unrecoverable scale, and ResPlan in units
that appear to be decimetres. That single fact constrains the entire
design downstream. Absolute measurements can be compared only within a
source, so anything that crosses sources must use either measurements
that are independent of scale, as the feature judge does, or a
representation in which scale has been normalised away, as the image judge
does.

\subsection{Conversion}

Each dataset loader yields a plan identifier and its geometry, and one
conversion routine performs the rest:

\begin{enumerate}
\item rotate the plan so that its dominant wall direction is horizontal,
since Swiss plans are drawn at arbitrary angles;
\item reduce each wall polygon to a centre line and a thickness;
\item apply the load-bearing threshold, then group the surviving wall
positions into grid lines, weighting each wall by its length and
expressing tolerances as a fraction of plan size;
\item place a column at every crossing contained by a load-bearing wall,
then add the columns the drawing marks;
\item hash the arrangement after translating it to the origin, so that
identical layouts drawn in different places collapse to one entry;
\item compute the geometric measurements.
\end{enumerate}

Only walls within ten degrees of an axis contribute to step 3. A wall at a
genuine angle is not straightened by any rotation, and forcing it into the
grouping would place a grid line where no wall exists. Such walls are
handled separately in the parser, as Section~\ref{sec:parser} describes.

\section{Deriving the labels}
\label{sec:labels}

An architectural drawing does not record a structural column layout. Since
accuracy cannot be measured without one, this work derives a label by a
stated procedure:

\begin{enumerate}
\item \textbf{Discard every wall thinner than 160\,mm.} This is the
minimum thickness for a two-hour fire rating, so thinner walls are treated
as room dividers rather than primary structure.
\item \textbf{Build the grid from the surviving walls only.} A structural
grid should follow structural elements; using every wall places grid lines
at room spacing rather than at structural spacing.
\item \textbf{Place a column at every grid crossing that lies inside one of
those walls.} The containment test uses each wall's start point, end point
and thickness, so a short wall claims only the crossings it actually
reaches. Testing the grid line alone would allow a stub wall at one end of
a building to claim crossings at the other.
\item \textbf{Add every column the drawing marks explicitly}, whether or
not it falls inside a wall.
\end{enumerate}

\subsection{What the result represents}

Steps 1 to 3 produce an \emph{equivalent-frame idealisation of a
wall-bearing structure}: a column layout that stands in for the wall
system it replaces. These buildings do not have a column grid waiting to
be recovered.

Badoux and Peter, surveying more than fifty Swiss buildings from original
construction drawings, report that in the dominant Swiss apartment type
the partition walls \emph{are} the structural bearing walls and that
frame-type structural systems are the exception. The corpus agrees: its
ratio of wall area to floor area is about 24 per cent, against roughly 15
per cent for shear-wall buildings and 8 per cent for concrete-frame
buildings in published surveys. Because nearly every wall carries load,
the 160\,mm threshold selects the \emph{primary} walls rather than
separating structure from non-structure.

Only step 4 is observed. Everything else is a modelling decision, and this
report is explicit about the distinction wherever a result depends on it.

\subsection{The threshold, and its sensitivity}

The 160\,mm value sits inside a narrow band of published figures. Eurocode
2 gives 140\,mm as the minimum for a reinforced load-bearing wall; Badoux
and Peter report 140\,mm as the thinnest concrete wall found in the Swiss
stock; the DISCS building inventory places its tenth percentile at 140 to
150\,mm across 102 Swiss buildings and its median structural wall at
200\,mm. The value used, 160\,mm, is the two-hour fire-rating minimum.

The measured thicknesses validate externally against that inventory: a
median of 190\,mm against DISCS's 200\,mm, and a seventy-fifth percentile
of 291\,mm against their 300\,mm. The additional mass below 150\,mm in
this corpus is the partition population that DISCS excludes by design,
since it inventories only load-bearing and insulating layers.

Because every downstream figure depends on this one number, the labels are
re-derived at 140, 160, 180 and 200\,mm. Table~\ref{tab:thresholdcorpus}
gives the effect at corpus level, and Section~\ref{sec:sensitivity}
carries it through to accuracy.

\begin{table}[H]
\centering
\caption[Effect of the load-bearing threshold on the corpus]{Effect of the
load-bearing threshold on the corpus, Swiss Dwellings. Raising the
threshold admits fewer walls, so fewer plans yield a usable label and the
labels that survive are sparser.}
\label{tab:thresholdcorpus}
\small
\begin{tabular}{rrr}
\toprule
Threshold (mm) & Plans with labels & Median columns \\
\midrule
140 & 8,830 & 52 \\
160 & 8,744 & 40 \\
180 & 8,587 & 30 \\
200 & 8,407 & 24 \\
\bottomrule
\end{tabular}
\end{table}

\subsection{What each source supports}

\textbf{Swiss Dwellings} is drawn in metres and stores walls as polygons,
so thickness is measurable and the threshold applies as written. It also
marks columns explicitly, supplying step 4. This is the only source where
the procedure works exactly as specified, and it is the source from which
all evaluation plans are drawn.

\textbf{CubiCasa5K} also stores walls as polygons, but in drawing pixels
with no recoverable real-world scale, so an absolute threshold in
millimetres is meaningless. It instead retains the thickest walls of each
plan by proportion. This is a genuine weakness and is why those plans
carry a different confidence label. Recovering scale from standard door
widths would resolve it and is left as further work.

\textbf{ResPlan} supports neither step 1 nor step 4: it declares one wall
depth for a whole plan, so there is no per-wall thickness to filter on, and
it marks no columns. Its plans remain on an earlier every-crossing
procedure, are labelled accordingly, and are excluded from every training
set. They remain useful for descriptive statistics and for exercising the
parser against a third drawing convention.

\subsection{Published artefacts}

Earlier revisions of this work retained only summary measurements and grid
lines, which meant the labels themselves were unavailable and neither the
judges nor the accuracy protocol could be constructed. The regeneration
script now publishes the ground-truth column coordinates of every labelled
plan and its load-bearing walls with thicknesses, and validates its output
against the previously committed tables before overwriting them. On the
most recent run the regenerated Swiss plans matched on all 8,260 shared
identifiers, with grid-line counts and regularity identical, and CubiCasa
matched on all 4,205.

\section{What the corpus looks like}
\label{sec:eda}

Two properties of the corpus govern what any learned component can
achieve.

\textbf{Building the grid from load-bearing walls only was the most
consequential change made to this pipeline.} Under an earlier procedure
that placed a column at every wall crossing, plans yielded a median of 84
to 130 candidate columns, where a structural grid for a dwelling is nearer
20 to 40, because every internal partition contributed a grid line. With
the grid built from load-bearing walls, the median falls to 40 for Swiss
and 11 for CubiCasa, and the proportion of plans with more than 100
columns falls from 71 to 9 per cent for Swiss and from 33 to 0.1 per cent
for CubiCasa.

\textbf{Real plans are irregular.} Regularity, defined so that it equals 1
when every bay is the same size and falls towards 0 as bay sizes diverge,
has a median of 0.34 to 0.47 across the three sources.
Figure~\ref{fig:regularity} shows the distributions. Published synthetic
training grids score far higher, and the consequence was measured
directly: a classifier trained only on synthetic arrangements scored 0.57
ROC AUC on real ones, where 0.5 is chance. That single figure is why every
learned component in this work trains on the real corpus.

\begin{figure}[H]
\centering
\includegraphics[width=0.85\textwidth]{03_regularity}
\caption[Regularity of real plans by source]{Regularity of real plans by
source. Values approach 1 when every structural bay is the same size. The
median lies between 0.34 and 0.47, far below the uniform grids used in
published synthetic training sets, which is why models trained on
synthetic layouts transfer poorly to real ones.}
\label{fig:regularity}
\end{figure}

Swiss Dwellings is the most regular source at 0.468, consistent with its
being multi-storey apartment blocks, which makes it the closest available
proxy for the commercial frames this pipeline ultimately targets.

\section{The parser}
\label{sec:parser}

The parser converts wall segments into a structural grid and a set of
candidate column positions. Each wall votes for a grid line: a wall
running vertically votes for a vertical line at its midpoint, weighted by
its length, and a horizontal wall votes correspondingly. Votes closer
together than a merge tolerance are combined into one line positioned at
the length-weighted average of its members. Candidate positions are the
crossings of the surviving lines. User constraints may exclude regions of
the plan, marking the crossings inside them unavailable.

Where the input records a thickness for each wall, only walls at or above
the load-bearing threshold may vote, and every crossing is tested for
whether it lies physically inside such a wall using that wall's own
thickness. This is the same test that produced the labels in
Section~\ref{sec:labels}.

That alignment is deliberate and it determines whether the evaluation
means anything. If the parser interpreted walls differently from the
labelling procedure, the correct arrangement could never appear among the
candidates, and every accuracy figure would measure the disagreement
between two parsers rather than the quality of any variant. Achieving
alignment required matching the labelling procedure in four respects:
midpoint votes, weighting by the wall's length component along its own
axis, a merge tolerance of 1.5 per cent of the plan extent, and silence
from walls shorter than 5 per cent of the plan. With those in place the
correct arrangement is reachable on all seven evaluation plans, as
Section~\ref{sec:coverage} reports.

Figure~\ref{fig:parsed} shows a parsed plan. Columns that the drawing
itself marks are treated as input rather than as part of the answer: the
parser records them and the generator retains them in every candidate.

\begin{figure}[H]
\centering
\includegraphics[width=0.8\textwidth]{04_parsed_dwelling}
\caption[A parsed plan]{A held-out dwelling after parsing. Grey lines are
wall centre lines; pale blue lines are the grid the load-bearing walls
imply. Dark markers are crossings that lie inside a load-bearing wall and
are therefore candidate column positions under the wall-based placement
rule; pale markers are crossings that do not.}
\label{fig:parsed}
\end{figure}

\section{The generator}
\label{sec:generator}

A candidate arrangement is a choice of grid lines together with a rule for
which crossings receive a column. The outermost line on each axis is
always retained, and any subset of the interior lines is permitted
provided every resulting bay width falls within the user's minimum and
maximum span. Crossings excluded by constraint are removed, and such a
removal is allowed only if the neighbouring columns can still span the
resulting gap.

Two placement rules are applied to each surviving set of lines. The
\emph{lattice} rule places a column at every available crossing, which is
the classic frame answer. The \emph{walls} rule places a column only at
crossings inside a load-bearing wall, which reproduces the procedure that
generated the labels. The second rule is available only where wall
thicknesses are known, and without it the correct arrangement would not be
among the candidates at all.

\subsection{Sampling large search spaces}

The number of line subsets grows exponentially with the number of interior
lines, as Figure~\ref{fig:searchspace} shows, so the space of candidate
pairs cannot always be enumerated in full. Where it exceeds the configured
limit the generator draws a random sample of the pairs using a fixed seed,
always including the pair that retains every line, which is where the
correct arrangement lies.

\begin{figure}[H]
\centering
\includegraphics[width=0.72\textwidth]{05_search_space}
\caption[Growth of the candidate space]{Growth of the candidate space with
the number of interior grid lines on one axis, logarithmic vertical scale.
Exhaustive enumeration is feasible for simple plans and infeasible for
heavily subdivided ones, which is why the generator samples beyond a
stated limit.}
\label{fig:searchspace}
\end{figure}

Sampling rather than truncating matters. Stopping at the limit would fill
the candidate set with near-identical dense grids and silently discard the
sparse end of the space; a sample keeps the set representative, and the
fixed seed makes it reproducible. Sampling reduces coverage without
restoring exhaustiveness, and that is stated as a limitation rather than
presented as a solution. At the default limit of 4,000 candidates the
seven evaluation plans produce between 512 and 4,000 candidates each.

\section{The score}
\label{sec:score}

Each candidate receives a single number, the penalty, computed from the
geometry of the arrangement and of the plan. Lower is better. The penalty
contains no member sizes, no load values, no prices and no units:
\begin{equation}
\text{penalty} = \sum_{i} w_i \, z(t_i),
\label{eq:penalty}
\end{equation}
where the $t_i$ are the five measures below, $z$ is a robust
standardisation within the candidate set of one plan, and the weights
$w_i$ are the function's only free parameters.

\subsection{The five measures}

Each measure is dimensionless and each increases as the arrangement gets
worse.

\begin{description}[style=unboxed,leftmargin=0pt]
\item[Density, $t_\mathrm{density}$.] Column count divided by floor area.
A proxy for the quantity of vertical structure and the number of
foundations, without pricing either. Replaces the material cost.

\item[Span demand, $t_\mathrm{span}$.] The sum over beam runs of
$s^2 b$, where $s$ is the span and $b$ the width of floor that run
supports, normalised by a footprint term. The bending a beam must resist
is proportional to $wL^2/8$ per unit of supported width; removing the
constant of proportionality, which is identical for every candidate,
leaves a quantity that still orders arrangements by structural demand
while requiring no load value and no steel section. This is the direct
replacement for the deleted member sizing, and the physics is preserved
exactly up to that constant. Spans are measured between the actual columns
on each grid line, so a sparse line with a long unsupported run is charged
for it.

\item[Irregularity, $t_\mathrm{irreg}$.] The variability of bay sizes in
each direction, summed. Regular grids give predictable load paths and
repeated details, and grid regularity is the standard early-stage measure
of layout quality.

\item[Off-wall fraction, $t_\mathrm{offwall}$.] The fraction of columns
not lying inside a wall. Columns away from walls intrude on the
architecture and require non-standard detailing. This is the only measure
that reads the architectural plan rather than the abstract grid, and it
uses the same containment test as the labelling procedure so that training
and use agree.

\item[Repetition, $t_\mathrm{rep}$.] The number of distinct bay sizes
divided by the number of bays, where sizes within 5 per cent of the mean
count as the same. Repeated bays mean repeated formwork and simpler
procurement, which is a buildability argument that survives the loss of a
section catalogue.
\end{description}

\subsection{Normalisation}

Each measure is converted to a robust standard score within the candidate
set of one plan: the value minus the median, divided by the interquartile
range, with a scaled median-absolute-deviation fallback when that range
degenerates. Scaling between the minimum and maximum was rejected because
a single extreme candidate would rescale every other score, making the
result depend on how many candidates happened to be generated. The median
and interquartile range are unaffected by values outside the central half
of the set.

\subsection{Obtaining the weights}
\label{sec:weights}

The weights are the only free parameters, and both of the following are
carried through the evaluation.

\textbf{Fixed weights} set every $w_i$ to one. Equal weighting is the
deliberate choice in the absence of any argued basis for preferring one
measure over another, and its justification is empirical rather than
rhetorical: Section~\ref{sec:sensitivity} perturbs each weight by
$\pm 50$ per cent and removes each measure in turn, and reports the
resulting change in the ranking. A weight vector that cannot be argued for
is acceptable precisely when the ranking can be shown to be insensitive to
it.

\textbf{Learned weights} are fitted by ranking. For each of 11,098 plans
the real derived arrangement is the positive example, and up to five
alternatives generated on that same plan's grid are the negatives: the
full set of crossings, and subsets of the grid lines under both placement
rules. This yields 52,724 paired comparisons. A logistic model without an
intercept is fitted to the differences between the normalised measures, so
that because the model is linear in those differences its coefficients
\emph{are} the weights.

Two other formulations were considered and rejected as unsound. Fitting a
model to reproduce the fixed score's output would be circular, since the
learned score would be a copy of the thing it is meant to be compared
against. Training on the judges' real-versus-altered examples would make
the learned score and the judge the same model applied twice, so that
comparing the pipeline with and without a judge would measure nothing.

One overlap is stated rather than argued away: the alternatives used here
and the altered examples used by the judges both include arrangements that
place columns at every crossing. The two sets of examples are otherwise
distinct.
